% =====================================================================
%  ÔN TẬP — BẢN DỄ
%
%  Khác bản đầy đủ (on-tap.tex) ở chỗ: mỗi chương chỉ giữ thứ tối thiểu
%  để làm được bài — nó là gì, công thức nào, một ví dụ số nhỏ. Bỏ hết
%  phần bẫy, phần ngoài giáo trình, và các bài dài.
%
%  Biên dịch: pdflatex -interaction=nonstopmode on-tap-de.tex   (HAI LẦN)
% =====================================================================
\documentclass[12pt]{article}

\usepackage[utf8]{inputenc}
\usepackage[T5]{fontenc}
\usepackage[vietnamese]{babel}
\usepackage[a4paper,margin=2.4cm]{geometry}
\usepackage{amsmath}
\usepackage{amssymb}
\usepackage{enumitem}
\usepackage{booktabs}
\usepackage{array}
\usepackage{hyperref}

\hypersetup{unicode=true, hidelinks, pdftitle={On tap ban de}}
\setlist{leftmargin=*, topsep=3pt, itemsep=2pt}
\emergencystretch=3em
\setlength{\parindent}{0pt}
\setlength{\parskip}{5pt}

\newcommand{\nho}{\textbf{Nhớ:}\ }

\title{\vspace{-1.5cm}\bfseries Ôn tập — bản dễ}
\author{Phân tích dữ liệu thông minh}
\date{}

\begin{document}
\maketitle
\vspace{-0.8cm}

\begin{center}
\fbox{\parbox{0.92\linewidth}{\centering
Mỗi phần chỉ có ba thứ: \textbf{nó là gì}, \textbf{công thức}, \textbf{một ví dụ}.\\
Cần chi tiết hơn thì xem bản đầy đủ \texttt{on-tap.tex}.
}}
\end{center}

\tableofcontents
\newpage

% =====================================================================
\section{Mô tả dữ liệu}

\subsection*{Nó là gì}

Tóm tắt một đống số thành vài con số dễ nói.

\begin{itemize}
\item \textbf{Trung bình} --- điểm cân bằng.
\item \textbf{Trung vị} --- số đứng giữa khi sắp thứ tự.
\item \textbf{Độ lệch chuẩn} --- dữ liệu tản ra bao xa quanh trung bình.
\end{itemize}

\subsection*{Công thức}

\[
\bar{x} = \frac{\text{tổng}}{n}
\qquad
s^2 = \frac{\sum (x_i-\bar{x})^2}{n-1}
\qquad
s = \sqrt{s^2}
\]

\nho mẫu số là $n-1$, không phải $n$.

\subsection*{Ví dụ}

Dữ liệu: $2,\ 4,\ 4,\ 5,\ 7,\ 8$. Trung bình $= 30/6 = 5$.

Lấy từng số trừ 5 rồi bình phương:
$9 + 1 + 1 + 0 + 4 + 9 = 24$.

\[
s^2 = \frac{24}{6-1} = 4{,}8
\qquad
s = \sqrt{4{,}8} = 2{,}19
\]

\subsection*{Hai hệ số tương quan}

Đo hai biến có đi cùng nhau không. Cả hai đều nằm trong khoảng từ $-1$ đến $1$.

\begin{center}
\begin{tabular}{@{}lll@{}}
\toprule
 & \textbf{Pearson} & \textbf{Spearman} \\
\midrule
Hỏi gì & hai biến có nằm trên một đường thẳng không & có cùng tăng cùng giảm không \\
Dùng khi & dữ liệu là số & dữ liệu là số hoặc là thứ hạng \\
\bottomrule
\end{tabular}
\end{center}

Gần $+1$ là cùng chiều mạnh, gần $-1$ là ngược chiều mạnh, gần $0$ là không liên quan.

\nho trung bình bị số quá lớn hoặc quá nhỏ kéo đi, trung vị thì không. Dữ liệu có
số bất thường (thu nhập, giá nhà) thì báo cáo trung vị.

% =====================================================================
\section{Khoảng tin cậy}

\subsection*{Nó là gì}

Ta chỉ đo được một mẫu, không đo được cả nước. Khoảng tin cậy là câu trả lời kèm
biên độ sai: \emph{``khoảng chừng từ đây tới đây''}.

\subsection*{Công thức}

\[
\text{con số đo được} \ \pm\ z^{*} \times SE
\]

Với tỉ lệ phần trăm:

\[
SE = \sqrt{\frac{\hat p\,(1-\hat p)}{n}}
\qquad\qquad
z^{*} = 1{,}96 \ \text{ cho độ tin cậy } 95\%
\]

\subsection*{Ví dụ}

Hỏi 850 người, có $67\%$ đồng ý. Vậy $\hat p = 0{,}67$ và $n = 850$.

\[
SE = \sqrt{\frac{0{,}67 \times 0{,}33}{850}} = 0{,}0161
\]
\[
0{,}67 \ \pm\ 1{,}96 \times 0{,}0161
\;=\; 0{,}67 \pm 0{,}032
\;=\; (0{,}64;\ 0{,}70)
\]

Đọc: \emph{``tin cậy 95\% rằng tỉ lệ thật nằm trong khoảng 64\% đến 70\%''}.

\nho mẫu càng lớn thì khoảng càng hẹp. Muốn khoảng hẹp đi một nửa thì phải lấy mẫu
lớn gấp bốn.

% =====================================================================
\section{Kiểm định giả thuyết}

\subsection*{Nó là gì}

Có một tuyên bố. Ta có dữ liệu. Hỏi: dữ liệu có đủ mạnh để bác bỏ tuyên bố đó không,
hay chênh lệch chỉ là may rủi.

\subsection*{Bốn bước}

\begin{enumerate}
\item Viết $H_0$ (tuyên bố cần thử) và $H_a$ (điều ta nghi ngờ).
\item Tính thống kê kiểm định ($z$ hoặc $t$).
\item So với ngưỡng, hoặc tính p-value.
\item \textbf{p-value nhỏ hơn $\alpha$ thì bác bỏ $H_0$.}
\end{enumerate}

\subsection*{Ba công thức đủ dùng cho phần lớn bài}

\begin{center}
\begin{tabular}{@{}ll@{}}
\toprule
\textbf{Đề cho gì} & \textbf{Dùng công thức} \\
\midrule
Một trung bình, có $s$ &
  $t = \dfrac{\bar x - \mu_0}{s/\sqrt n}$ \quad với $df = n-1$ \\[10pt]
Một tỉ lệ &
  $z = \dfrac{\hat p - p_0}{\sqrt{p_0(1-p_0)/n}}$ \\[10pt]
Hai trung bình, mẫu lớn &
  $z = \dfrac{\bar x_1 - \bar x_2}{\sqrt{s_1^2/n_1 + s_2^2/n_2}}$ \\
\bottomrule
\end{tabular}
\end{center}

\subsection*{Ngưỡng nào}

\begin{center}
\begin{tabular}{@{}lll@{}}
\toprule
\textbf{Đề hỏi} & \textbf{Kiểu} & \textbf{Ngưỡng khi } $\alpha = 0{,}05$ \\
\midrule
``có khác không'' & hai phía & $|z| > 1{,}96$ \\
``có tăng không'' & một phía & $z > 1{,}645$ \\
``có giảm không'' & một phía & $z < -1{,}645$ \\
\bottomrule
\end{tabular}
\end{center}

\subsection*{Ví dụ}

Giám đốc nói lương trung bình là 380. Lấy mẫu $n = 36$ người thấy $\bar x = 350$ và
$s = 40$. Kiểm ở $\alpha = 0{,}05$.

$H_0$: lương $= 380$. \quad $H_a$: lương $\neq 380$ (hai phía).

\[
t = \frac{350 - 380}{40/\sqrt{36}} = \frac{-30}{6{,}667} = -4{,}50
\]

Ngưỡng là $2{,}030$. Vì $4{,}50 > 2{,}030$ nên \textbf{bác bỏ $H_0$} --- lương thật
sự thấp hơn 380.

\nho không bác bỏ được thì viết \emph{``chưa đủ bằng chứng''}, đừng viết
\emph{``chấp nhận $H_0$''}.

% =====================================================================
\section{Mô hình đồ thị}

\subsection*{Nó là gì}

Vẽ mũi tên giữa các biến để biết cái nào ảnh hưởng cái nào, rồi từ đó biết khi
phân tích thì \emph{phải giữ cố định biến nào}.

\subsection*{Ba hình cần thuộc}

\begin{center}
\begin{tabular}{@{}llp{5.6cm}@{}}
\toprule
\textbf{Hình} & \textbf{Tên} & \textbf{Có kiểm soát không} \\
\midrule
$X \leftarrow Z \to Y$ & gây nhiễu & \textbf{CÓ}. Không kiểm soát thì kết quả sai. \\
$X \to Z \to Y$ & trung gian & \textbf{KHÔNG}. Kiểm soát là xoá mất cái cần đo. \\
$X \to Z \leftarrow Y$ & hội tụ & \textbf{KHÔNG}. Kiểm soát là tự tạo ra liên hệ giả. \\
\bottomrule
\end{tabular}
\end{center}

Cách nhận ra: nhìn chiều mũi tên ở biến giữa. Hai mũi tên \emph{đâm vào} nó thì đó
là hội tụ.

\subsection*{Ví dụ}

Muốn biết \emph{học thêm} có làm \emph{điểm} cao lên không. Nhưng \emph{gia đình
khá giả} vừa làm trẻ dễ đi học thêm hơn, vừa làm điểm cao hơn.

\[
\text{học thêm} \ \leftarrow\ \text{gia đình khá giả} \ \to\ \text{điểm}
\]

Đây là \textbf{gây nhiễu}. Không giữ cố định mức khá giả thì ta sẽ tưởng học thêm
hiệu quả hơn thực tế.

\nho chỉ kiểm soát biến \emph{gây nhiễu}. Nhét hết mọi biến vào mô hình cho chắc là
sai.

% =====================================================================
\section{Phân tích nhân quả}

\subsection*{Nó là gì}

Phân biệt \emph{đi cùng nhau} với \emph{gây ra}.

\subsection*{Ý chính}

Với mỗi người ta chỉ thấy được \emph{một} kết cục: hoặc là có uống thuốc, hoặc là
không. Cái còn lại vĩnh viễn không quan sát được. Nên không đo được tác động lên
từng cá nhân, chỉ đo được \textbf{trung bình cả nhóm}:

\[
ATE = \text{trung bình nhóm có} \;-\; \text{trung bình nhóm không}
\]

Nhưng phép trừ này \textbf{chỉ đúng khi hai nhóm được chia ngẫu nhiên}. Chia ngẫu
nhiên thì hai nhóm giống nhau ở mọi mặt khác, nên chênh lệch còn lại đúng là do
can thiệp.

\subsection*{Ví dụ ngược đời}

Người tập gym khoẻ hơn người không tập. Có phải gym làm khoẻ không?

Chưa chắc --- \emph{người vốn đã khoẻ mới hay đi gym}. Hai nhóm khác nhau ngay từ
đầu, nên phần chênh lệch không hoàn toàn là công của gym.

\subsection*{Nghịch lý Simpson}

Cách chữa A tốt hơn B ở \emph{cả} nhóm sỏi nhỏ ($93\%$ so với $87\%$) lẫn nhóm sỏi
lớn ($73\%$ so với $69\%$). Nhưng gộp chung lại thì A chỉ được $78\%$ còn B được
$83\%$ --- B thắng.

Lý do: bác sĩ dành cách A cho ca nặng. Kích thước sỏi là \textbf{biến gây nhiễu}.

\nho gặp trường hợp này thì tin con số \emph{theo từng nhóm}, đừng tin con số gộp.

% =====================================================================
\section{Hồi quy tuyến tính}

\subsection*{Nó là gì}

Vẽ một đường thẳng đi xuyên qua đám điểm, để từ $x$ đoán ra $y$.

\subsection*{Công thức}

\[
\hat y = b_0 + b_1 x
\qquad
b_1 = \frac{S_{xy}}{S_{xx}}
\qquad
b_0 = \bar y - b_1 \bar x
\]

$b_1$ là \textbf{độ dốc}: $x$ tăng 1 thì $y$ đổi bấy nhiêu.

$R^2$ cho biết đường thẳng giải thích được bao nhiêu phần trăm biến động của $y$.
$R^2 = 0{,}73$ nghĩa là $73\%$.

\subsection*{Ví dụ}

\begin{center}
\begin{tabular}{@{}lccccc@{}}
\toprule
$x$ & 1 & 2 & 3 & 4 & 5 \\
$y$ & 2 & 4 & 5 & 4 & 9 \\
\bottomrule
\end{tabular}
\end{center}

Tính được $\bar x = 3$, $\bar y = 4{,}8$, $S_{xy} = 14$, $S_{xx} = 10$.

\[
b_1 = \frac{14}{10} = 1{,}4
\qquad
b_0 = 4{,}8 - 1{,}4 \times 3 = 0{,}6
\]

Vậy $\hat y = 0{,}6 + 1{,}4x$. Đọc: $x$ tăng 1 thì $y$ tăng khoảng $1{,}4$.

\textbf{Cách tự kiểm tra:} thay $x = \bar x = 3$ vào phải ra đúng $\bar y = 4{,}8$.
Thử: $0{,}6 + 1{,}4 \times 3 = 4{,}8$. Đúng.

\nho dữ liệu quan sát thì viết \emph{``đi kèm''}, đừng viết \emph{``làm cho''}.

% =====================================================================
\section{Hồi quy đa biến và logistic}

\subsection*{Đa biến --- nhiều $x$ cùng lúc}

\[
\hat y = b_0 + b_1 x_1 + b_2 x_2 + \dots
\]

Đọc $b_1$: $x_1$ tăng 1 đơn vị thì $y$ đổi bấy nhiêu, \textbf{khi các biến còn lại
giữ nguyên}. Cụm in đậm bắt buộc phải viết.

\subsection*{Logistic --- khi $y$ chỉ có hai giá trị}

Dùng khi $y$ là có/không, đậu/rớt, sống/chết. Không dùng đường thẳng được vì đường
thẳng có thể cho ra xác suất âm hoặc lớn hơn 1.

Ba bước, làm đúng thứ tự:

\begin{enumerate}
\item Thay số vào để được \textbf{logit}.
\item $\text{odds} = e^{\text{logit}}$.
\item $\text{xác suất} = \dfrac{\text{odds}}{1+\text{odds}}$.
\end{enumerate}

\subsection*{Ví dụ}

Cho $\text{logit} = 1{,}8185 - 0{,}0665 \times \text{tuổi}$. Tính xác suất sống sót
của người 25 tuổi.

\begin{enumerate}
\item $\text{logit} = 1{,}8185 - 0{,}0665 \times 25 = 0{,}156$
\item $\text{odds} = e^{0{,}156} = 1{,}169$
\item $\text{xác suất} = \dfrac{1{,}169}{1 + 1{,}169} = 0{,}539$
\end{enumerate}

Vậy khoảng $54\%$.

\nho $e^{\beta}$ là \emph{tỉ số odds}, không phải tỉ số xác suất. Hệ số $\beta = 0$
nghĩa là không ảnh hưởng, vì $e^0 = 1$.

% =====================================================================
\section{Monte Carlo}

\subsection*{Nó là gì}

Bài toán khó tính bằng công thức thì \emph{bốc số ngẫu nhiên rất nhiều lần rồi đếm}.

\subsection*{Ví dụ kinh điển --- tính số $\pi$}

Vẽ một hình vuông và đường tròn nội tiếp. Ném ngẫu nhiên $n$ điểm vào hình vuông,
đếm được $k$ điểm rơi vào trong tròn. Khi đó:

\[
\pi \approx 4 \times \frac{k}{n}
\]

Càng ném nhiều càng chính xác.

\subsection*{Sai số}

\[
SE = \frac{\sigma}{\sqrt{n}}
\]

\nho muốn sai số giảm một nửa thì phải chạy \textbf{gấp bốn lần} số vòng.

\subsection*{Bootstrap}

Có đúng một mẫu gồm $n$ món, muốn biết kết quả của mình chắc chắn tới đâu:

\begin{enumerate}
\item Bốc \textbf{có hoàn lại} đúng $n$ món (bốc xong bỏ lại vào, nên có món trúng
      hai lần, có món không trúng lần nào).
\item Tính lại con số cần tính.
\item Lặp vài nghìn lần.
\item Độ tản của mấy nghìn kết quả đó chính là sai số.
\end{enumerate}

% =====================================================================
\section{Bảng tra}

\begin{center}
\begin{tabular}{@{}lcccc@{}}
\toprule
\textbf{Độ tin cậy} & 90\% & 95\% & 98\% & 99\% \\
\midrule
$z^{*}$ hai phía & 1,645 & \textbf{1,96} & 2,33 & 2,576 \\
$z$ một phía & 1,282 & \textbf{1,645} & 2,054 & 2,326 \\
\bottomrule
\end{tabular}
\end{center}

Thuộc ba số \textbf{1,96} --- \textbf{1,645} --- \textbf{2,576} là gỡ được phần lớn bài.

\subsection*{Năm câu đừng viết sai}

\begin{enumerate}
\item p-value \emph{không} phải xác suất $H_0$ đúng.
\item Viết \emph{``chưa đủ bằng chứng để bác bỏ''}, đừng viết \emph{``chấp nhận $H_0$''}.
\item Phương sai chia $n-1$; hồi quy dùng $df = n-2$.
\item Dữ liệu quan sát thì viết \emph{``đi kèm''}, đừng viết \emph{``làm cho''}.
\item Chỉ kiểm soát biến \emph{gây nhiễu}, không kiểm soát biến trung gian và hội tụ.
\end{enumerate}

\end{document}
